\section{Preliminaries and Notation}
\label{sec:prelim}

\paragraph{Basic notation.}
For an integer \(m\ge 1\), write \([m]:=\{1,\dots,m\}\).
We work over \(\F_2=\{0,1\}\) with addition \(\oplus\) (equivalently, addition modulo \(2\)).
For \(x\in\{0,1\}^m\), let \(\wt(x)\) denote its Hamming weight and write \(p(x):=\wt(x)/m\).

\paragraph{Probability and expectation.}
All probabilities and expectations are taken with respect to the randomness explicitly specified in context.
We write \(\PP[\cdot]\) for probability and \(\EE[\cdot]\) for expectation. Conditional expectation is written
\(\EE[\cdot\mid \mathcal{F}]\) for a \(\sigma\)-algebra \(\mathcal{F}\).

\paragraph{Uniform distribution.}
For a finite set \(\mathcal{S}\), we write \(X\sim \Unif(\mathcal{S})\) to denote that \(X\) is uniform over \(\mathcal{S}\).
Common examples include:
\[
I\sim \Unif([m]) \quad\text{(a uniform index)},\qquad
\Pi\sim \Unif(S_m) \quad\text{(a uniform permutation)}.
\]
For sampling a uniform \(d\)-subset of \([m]\) (without replacement) we write
\[
S\sim \Unif\!\Big(\binom{[m]}{d}\Big).
\]

\paragraph{Permutations.}
Let \(S_m\) denote the symmetric group on \([m]\).
For \(\Pi\in S_m\) and \(x\in\{0,1\}^m\), we write \(\Pi(x)\) for the coordinate permutation
\[
(\Pi(x))_i := x_{\Pi^{-1}(i)}.
\]

\paragraph{Standard distributions.}
We use the following classical distributions.

\subsubsection*{Beta distribution}
For parameters \(a,b>0\), the Beta distribution \(\Beta(a,b)\) is supported on \([0,1]\) with density
\[
f(s) = \frac{1}{B(a,b)}\, s^{a-1}(1-s)^{b-1},
\]
where \(B(a,b)=\Gamma(a)\Gamma(b)/\Gamma(a+b)\) is the Beta function.

\subsubsection*{Dirichlet distribution}
For parameters \(\alpha_1,\dots,\alpha_k>0\), the Dirichlet distribution \(\Dirichlet(\alpha_1,\dots,\alpha_k)\)
is supported on the simplex
\[
\Delta_k := \Bigl\{(z_1,\dots,z_k)\in\R_{\ge 0}^k:\ \sum_{i=1}^k z_i = 1\Bigr\},
\]
with density proportional to \(\prod_{i=1}^k z_i^{\alpha_i-1}\).
The special case \(\Dirichlet(1,\dots,1)\) is uniform on \(\Delta_k\).

\subsubsection*{Multinomial distribution}
For \(N\in\Z_{\ge 0}\) and \(p=(p_1,\dots,p_k)\in\Delta_k\), the multinomial distribution \(\Multinomial(N;p)\)
is a distribution on \((K_1,\dots,K_k)\in\Z_{\ge 0}^k\) satisfying \(\sum_i K_i=N\), with probability mass function
\[
\PP[(K_1,\dots,K_k)=k]
=
\frac{N!}{\prod_{i=1}^k k_i!}\prod_{i=1}^k p_i^{k_i}.
\]
For any subset \(I\subseteq[k]\), the aggregate \(\sum_{i\in I}K_i\) is binomial:
\[
\sum_{i\in I}K_i\ \big|\ p \sim \Binomial\Bigl(N,\ \sum_{i\in I}p_i\Bigr).
\]

\paragraph{Dirichlet subset-sum marginal (Beta marginal).}
If \((Z_1,\dots,Z_k)\sim \Dirichlet(\alpha_1,\dots,\alpha_k)\) and \(I\subseteq[k]\), then
\[
\sum_{i\in I} Z_i \sim \Beta\Bigl(\sum_{i\in I}\alpha_i,\ \sum_{i\notin I}\alpha_i\Bigr).
\]
In particular, for \((Z_1,\dots,Z_k)\sim \Dirichlet(1,\dots,1)\) and \(|I|=r\),
\[
\sum_{i\in I} Z_i \sim \Beta(r,k-r).
\]